\subsection{What is a multi-agent system?}
\label{sec:what-is-a-mas}
\slidesrc{Intro to agents \& MAS: 12--17}

A single agent is rarely enough: interesting problems require several agents, each with
its own goals and abilities, that interact. Such a system is a \textbf{multi-agent system}
(MAS). Slide 12 shows the classic picture, redrawn in Figure~\ref{fig:mas-structure}.

\begin{center}
\begin{tikzpicture}[font=\small, >=Stealth,
  ag/.style={ellipse, draw=nbgreen!60!black, fill=nbgreen!35, thick, minimum width=1.4cm, minimum height=0.75cm, font=\small\bfseries}]
  % environment
  \draw[thick, nbgray, fill=nbgray!15] (5.2,-2.6) ellipse (5.0cm and 0.9cm);
  \node[font=\small\bfseries, nbgray!70!black] at (5.2,-3.85) {Environment};
  \fill[nbblue!25, draw=nbblue] (2.4,-2.55) ellipse (0.9cm and 0.3cm);
  \fill[nbblue!25, draw=nbblue] (4.1,-2.55) ellipse (0.9cm and 0.3cm);
  \fill[nbblue!25, draw=nbblue] (7.6,-2.45) ellipse (0.8cm and 0.3cm);
  % organizations
  \draw[thick, nbgray] (2.9,0.6) ellipse (2.7cm and 1.35cm);
  \draw[thick, nbgray] (7.1,0.75) ellipse (2.9cm and 1.45cm);
  \node[font=\small\bfseries] at (0.4,2.1) {Organisation};
  % agents
  \node[ag] (a1) at (1.1,0.2) {Agent};
  \node[ag] (a2) at (3.0,1.2) {Agent};
  \node[ag] (a3) at (5.0,0.6) {Agent};
  \node[ag] (a4) at (6.9,1.6) {Agent};
  \node[ag] (a5) at (9.0,1.0) {Agent};
  \node[ag] (a6) at (7.4,0.0) {Agent};
  \draw[<->, thick] (a1) -- (a2) node[midway, below right=-2pt, font=\scriptsize, align=left]{inter-agent\\communication};
  \draw[<->, thick] (a2) -- (a3);
  \draw[<->, thick] (a3) -- (a4);
  \draw[<->, thick] (a4) -- (a5);
  \draw[<->, thick] (a4) -- (a6);
  \draw[<->, thick] (a3) -- (a6);
  % access
  \draw[<->, ultra thick, nborange] (1.4,-0.3) -- (2.3,-2.2);
  \draw[<->, ultra thick, nborange] (4.8,0.15) -- (4.2,-2.2);
  \draw[<->, ultra thick, nborange] (7.5,-0.45) -- (7.6,-2.1);
  \node[font=\scriptsize, align=left] at (9.2,-1.4) {access to the\\environment\\(sphere of influence)};
\end{tikzpicture}
\captionof{figure}{Structure of a MAS (slide 12). Agents (green) communicate with each other,
are grouped in organisations (grey ellipses) and each one can perceive and act on its own
portion of the shared environment (blue regions, orange arrows). Spheres of influence may
overlap: two agents may act on the same part of the environment.}
\label{fig:mas-structure}
\end{center}

\begin{nbdef}{Characteristics of a MAS (slides 14--17)}
In a multi-agent system:
\begin{enumerate}[nosep]
  \item each agent has \textbf{incomplete information}, or capabilities for solving the
        problem, thus each agent has a \textbf{limited viewpoint};
  \item there is \textbf{no global system control};
  \item data is \textbf{decentralised};
  \item computation is \textbf{asynchronous}.
\end{enumerate}
\end{nbdef}

\begin{nbexample}{a smart-city ride-sharing MAS}
\notinslides
Passenger agents (one per user app), driver agents (one per car) and a pricing agent.
\begin{enumerate}[nosep]
  \item \emph{Limited viewpoint}: a driver agent knows its own position and route, not the
        positions of the other cars or the requests of all passengers.
  \item \emph{No global control}: no central program decides every move; drivers accept or
        refuse rides.
  \item \emph{Decentralised data}: each phone stores its own user's preferences and history.
  \item \emph{Asynchronous computation}: requests and answers arrive at any time; nobody
        waits for a global ``clock tick''.
\end{enumerate}
The same four features describe e-commerce marketplaces, power grids with home batteries,
fleets of warehouse robots and swarms of drones.
\end{nbexample}

\begin{nbintuition}[Why ``no global control'' is a feature, not a bug]
A central controller is a single point of failure and a bottleneck, and in an \emph{open}
system (agents of different owners, joining and leaving at any time) it simply cannot exist:
nobody owns all the agents. MAS replace central control by \emph{interaction}: negotiation,
coordination, competition. This is why agent \emph{communication}
(Section~\ref{sec:agent-communication}) is a central topic of the course.
\end{nbintuition}
